\documentclass[10pt,a3paper,landscape]{article}
\usepackage[margin=6mm]{geometry}
\usepackage[T1]{fontenc}
\usepackage[utf8]{inputenc}
\usepackage[french]{babel}
\usepackage{lmodern}
\usepackage{microtype}
\makeatletter\def\input@path{{../../../Style/}}\makeatother
\NeedsTeXFormat{LaTeX2e}
\ProvidesPackage{TLPosterCarteMentale}[2026/08/03 Posters et cartes mentales SI]

\RequirePackage{tikz}
\RequirePackage[most]{tcolorbox}
\RequirePackage{xcolor}
\RequirePackage{amsmath,amssymb}
\RequirePackage{graphicx}
\usetikzlibrary{positioning,arrows.meta,calc,shapes.geometric,mindmap,backgrounds,fit}

\definecolor{TLPosterBlue}{RGB}{24,86,141}
\definecolor{TLPosterLight}{RGB}{234,243,250}
\definecolor{TLPosterAccent}{RGB}{232,126,34}
\definecolor{TLPosterGreen}{RGB}{52,152,102}
\definecolor{TLPosterRed}{RGB}{192,57,43}
\definecolor{TLPosterGray}{RGB}{80,90,100}

\newcommand{\TLPosterTitre}[2]{%
  \begin{tcolorbox}[enhanced,boxrule=0pt,arc=0pt,colback=TLPosterBlue,left=6mm,right=6mm,top=4mm,bottom=4mm]
    {\color{white}\bfseries\LARGE #1}\hfill{\color{white}\large #2}
  \end{tcolorbox}%
}

\newtcolorbox{TLPosterBloc}[2][]{
  enhanced,breakable,
  colback=TLPosterLight,
  colframe=TLPosterBlue,
  colbacktitle=TLPosterBlue,
  coltitle=white,
  fonttitle=\bfseries,
  title={#2},
  boxrule=0.8pt,
  arc=2mm,
  left=3mm,right=3mm,top=2mm,bottom=2mm,
  #1
}

\newtcolorbox{TLPosterFormule}[1][]{
  enhanced,
  colback=white,
  colframe=TLPosterAccent,
  boxrule=1pt,
  arc=2mm,
  fontupper=\bfseries,
  halign=center,
  #1
}

\newcommand{\TLPosterMethode}[1]{%
  \begin{tcolorbox}[enhanced,colback=TLPosterGreen!8,colframe=TLPosterGreen,title=Méthode,fonttitle=\bfseries]
  #1
  \end{tcolorbox}%
}

\newcommand{\TLPosterAttention}[1]{%
  \begin{tcolorbox}[enhanced,colback=TLPosterRed!6,colframe=TLPosterRed,title=Point de vigilance,fonttitle=\bfseries]
  #1
  \end{tcolorbox}%
}

\tikzset{
  TLmindroot/.style={concept,concept color=TLPosterBlue,text=white,font=\bfseries\large,minimum size=34mm},
  TLmindlevel1/.style={level 1 concept/.append style={font=\bfseries,minimum size=23mm,sibling angle=45}},
  TLmindlevel2/.style={level 2 concept/.append style={font=\small,minimum size=17mm}},
  TLarrow/.style={-{Stealth[length=3mm]},thick,draw=TLPosterBlue},
  TLbox/.style={rounded corners=2mm,draw=TLPosterBlue,fill=TLPosterLight,align=center,inner sep=3mm}
}

\newenvironment{TLCarteMentale}[1]{%
  \begin{tikzpicture}[mindmap,grow cyclic,concept color=TLPosterBlue,every node/.style={concept},TLmindlevel1,TLmindlevel2]
  \node[TLmindroot]{#1}
}{;
  \end{tikzpicture}
}

\newcommand{\TLBranche}[3][]{child[concept color=#2]{node[#1]{#3}}}

\endinput

\pagestyle{empty}
\renewcommand{\familydefault}{\sfdefault}
\begin{document}
\begin{tikzpicture}
\TLMapBase{Logique et SED}{Identifier le cas $\rightarrow$ appliquer la méthode adaptée $\rightarrow$ vérifier tous les scénarios}
\TLMapBranchTL{À connaître}{\textbullet\ combinatoire : pas de mémoire\\[1mm]\textbullet\ séquentiel : état mémorisé\\[1mm]\textbullet\ NON, ET, OU, XOR, De Morgan\\[1mm]\textbullet\ Karnaugh et formes canoniques\\[1mm]\textbullet\ état / transition / événement / garde / effet\\[1mm]\textbullet\ codeurs A/B/Z et Gray\\[1mm]\textbullet\ parité, CRC, débit, délai, BER, OSI}
\TLMapBranchTR{Méthode combinatoire}{1. Définir entrées/sorties et états actifs.\\[1mm]2. Construire la table de vérité.\\[1mm]3. Écrire les mintermes des sorties à 1.\\[1mm]4. Simplifier par Boole ou Karnaugh.\\[1mm]5. Karnaugh : groupes de puissances de 2, les plus grands possibles ; bords adjacents ; $X$ utilisables.\\[1mm]6. Réaliser et tester toutes les combinaisons utiles.}
\TLMapBranchML{Méthode SED / états}{1. Recenser les états stables.\\[1mm]2. Définir les activités entry/do/exit.\\[1mm]3. Tracer les transitions : événement [garde] / effet.\\[1mm]4. Ajouter initial, choix, after(T), fork/join, historique si nécessaire.\\[1mm]5. Vérifier complétude, exclusivité des gardes, sécurité, accessibilité et synchronisation des régions parallèles.}
\TLMapBranchMR{Méthode codeur}{\textbf{Incrémental :} $N$ périodes/tour ; sur 4 fronts $N_c=4N$ ; $\Delta\theta=2\pi/N_c$. A avant B / B avant A donne le sens.\\[1mm]\textbf{Position :} compteur signé.\\[1mm]\textbf{Vitesse :} comptage sur fenêtre à haute vitesse ; période entre fronts à basse vitesse.\\[1mm]\textbf{Absolu :} $n$ pistes $\Rightarrow2^n$ positions ; Gray change un seul bit à la fois.}
\TLMapBranchBL{Méthode réseau / trame}{\textbf{Parité :} compter les 1 et appliquer paire/impaire.\\[1mm]\textbf{CRC :} ajouter $r$ zéros $\rightarrow$ division modulo 2 $\rightarrow$ insérer reste $\rightarrow$ vérifier reste nul.\\[1mm]\textbf{Choix réseau :} débit utile, délai/gigue, déterminisme, BER, support, topologie, simplex/duplex, série/parallèle, synchronisation et acquittement.}
\TLMapBranchBR{Décider / vérifier}{\textbf{Sans mémoire :} table de vérité.\\[1mm]\textbf{Avec histoire du système :} graphe d'états.\\[1mm]\textbf{Mesure numérique :} résolution + décodage.\\[1mm]\textbf{Communication :} structure de trame + contrôle d'erreur + contraintes temporelles.\\[1mm]\textbf{Toujours :} états impossibles, défauts, priorités, transitions invalides, ordre chronologique.}
\TLMapRibbon{Combinatoire $\neq$ séquentiel : le premier se simplifie par équations ; le second se valide par scénarios et transitions.}
\end{tikzpicture}
\end{document}
